\section{Related work}\label{sec:related}

The reference point for the modern literature on algorithmic collusion is \citet{calvano2020}; \citet{bichler2025} survey the literature that followed. Two or more independent tabular Q-learners set prices in a repeated Bertrand game with differentiated products, each conditioning its price on the previous period's price profile. They converge to prices well above the static Nash level, and the authors summarise the outcome with a profit index that places realised profit between the static Nash profit (zero) and the joint-monopoly profit (one). Their central piece of evidence is a deviation test: a firm is forced to cut its price for one period, and the rivals respond with a temporary price cut followed by a return to the old level. That punish-then-forgive pattern is the ground on which they defend the word ``collusion'': it has the shape of the reward-punishment strategies on which the folk theorem's equilibria rest. \citet{klein2021} obtains supra-competitive prices under sequential (alternating-move) pricing and reports both constant-price outcomes and cyclical ones that resemble Edgeworth cycles. Earlier, \citet{waltman2008} found that Q-learners in a Cournot game produce less than the Cournot--Nash quantity, and that they do so even without memory, so the outcome cannot rest on a punishment strategy; that observation anticipates the memory-0 result reported below. \citet{asker2022} make the point explicitly: whether learners end up at competitive or supra-competitive prices depends on the learning protocol, in their case whether an agent updates only the action it played or all actions at once, and not only on the incentives of the game. The present paper uses a price index $\Dl$ with cost as the zero (Section~\ref{sec:model}). Under uniform pricing with inelastic demand industry profit is $(p-10)\times 100$, so $\Dl$ equals an industry-profit index with cost as the zero; it is not the per-firm, Nash-zeroed index of \citet{calvano2020}, since firms earn unequal profits at every rest point reported below, and under pay-as-bid the price and profit forms diverge. It takes the deviation test and the memory-1 learner from this line and applies them in a market whose payoff structure differs from differentiated-product Bertrand competition: a homogeneous good, capacity constraints, inelastic demand and uniform-price clearing. One consequence, that a lone undercut is paid the rivals' price and leaves the clearing price unchanged, is what makes the unclaimed temptation large and the forced deviation profitable here (Section~\ref{sec:discussion}); whether the negative result depends on it was not tested, since no market in which a lone undercut moves the price was run.

A second line asks whether seemingly collusive outcomes are collusion at all. \citet{abada2023} let independent learners operate storage in an electricity setting with imperfect monitoring and observe prices and profits that look collusive. They attribute the outcome to imperfect exploration rather than to a learned strategy, and show that interventions during learning, or a different learning arrangement, can move the market toward the competitive outcome. \citet{abada2024} extend the argument across a range of algorithms and find that greater sophistication does not by itself raise profits; supra-competitive profits can be a mistake of the learning process. Of the papers above, this one is closest in argument to \citet{abada2023} and \citet{abada2024}: it agrees that supra-competitive prices from Q-learners need not indicate a learned reward-punishment scheme. For the memory-0 learner it reaches a verdict of the same broad kind, that the price comes from how $\varepsilon$-greedy learners explore, but not the same mechanism: the price is not an accident of exploration ending too soon, since it survives ten-fold slower decay and a restart from $\varepsilon_0 = 0.5$ (Sections~\ref{sec:beta} and~\ref{sec:reinject}). For the learner with memory it finds that in 69\% of temptations the table's own continuation value after deviating is low enough that one fresh update would not change the decision, which is what a learned continuation looks like (Section~\ref{sec:residual}); the deviation estimates themselves are as stale as the memory-0 learner's, and the continuation value comes from the same table and may be stale too, so whether this falls outside the ``mistake'' reading is left open (Section~\ref{sec:discussion}). It also works in a different market, a non-pivotal uniform-price auction with three symmetric firms. It runs a memory-0 control, a learner that cannot condition on history and therefore cannot carry a threat, alongside the memory-1 learner and applies the same indicator battery to both, and it adds a diagnostic that reads the Q-table itself, a one-step Bellman residual on the unplayed deviation, and asks whether a single fresh update would change the learner's decision. To our knowledge neither paper applies the forced-deviation test, the ablation and the temptation test to a memoryless learner as a calibration control alongside a learner with memory, and neither reports a Bellman residual on the unplayed deviation; \citet{waltman2008} report the memoryless outcome but do not use it to calibrate indicators. % TODO-FACT: check Abada et al.\ (2024), which compares algorithms of varying sophistication and may include stateless learners, before keeping this sentence.

The third line is auctions. \citet{banchio2022} study repeated first- and second-price auctions with Q-learning bidders. In first-price auctions the bidders settle at bids well below their values while in second-price auctions they do not. They trace the difference to the one-increment overbidding incentive of a first-price auction, which lets bidders re-coordinate on low bids after experimentation. Their bidders carry no state, so the outcome cannot rest on a punishment strategy; the mechanism they describe is a property of the learning dynamic on a discrete bid grid, and it is the closest in the literature to the one proposed here for the memory-0 learner. % TODO-FACT: confirm their Q-learners are stateless.
On the market side, \citet{vonderfehr1993} analyse the uniform-price auction of the early British pool as a game between capacity-constrained generators and show that the equilibrium depends on whether a firm is pivotal: when demand can be met without it, competition drives the price to cost, and when it cannot, the price rises toward the cap. \citet{fabra2006} compare uniform-price and discriminatory (pay-as-bid) auctions under the same capacity structure and characterise the equilibria of each, in pure strategies under uniform pricing and in mixed strategies for the discriminatory auction when a firm is pivotal. Their low-demand case, in which no firm is pivotal, gives the same equilibrium outcome as the market here, price at cost; it is not the same payoff structure, because with two firms the low bidder serves all of demand and sets the price, so a lone undercut moves the price there and does not here. The static benchmark in this paper is not taken from them but enumerated directly: the pure Nash equilibria of the one-shot game on the 15-rung grid (Section~\ref{sec:model}), which reduce to their competitive equilibrium up to grid coarseness. Closest in setting is \citet{seredynski2026}, who model strategic bidding in an algorithmic electricity market as a repeated game with imperfect public monitoring, simulate it with multi-agent reinforcement learning, and report cases in which the agents sustain supra-competitive outcomes that satisfy the usual collusion indicators. This paper adopts the market structure of that study reduced to a single node, with no network, no congestion and one demand level, and replaces their agents with tabular Q-learners whose tables can be inspected directly; the ``replication'' of Section~\ref{sec:intro} is of the indicator battery of \citet{calvano2020}, applied to a new market, not of the learners or results of \citet{seredynski2026}, which are not reproduced.

\citet{maskin1988} give the theory of Edgeworth cycles: with alternating moves, prices are undercut step by step until a firm finds it more profitable to raise its price and restart the cycle. The paper uses that shape as an analogy for the pattern it argues the memory-0 learners trace under exploration (Section~\ref{sec:occupation}); moves here are simultaneous, no equilibrium is claimed, and the argument is verbal. \citet{klein2021} finds such cycles in the greedy play of Q-learners with memory under alternating moves; here the memory-0 learner freezes to a fixed point once exploration is off (Section~\ref{sec:results}), and the cycle is inferred, not observed: what the constant-$\varepsilon$ runs record is a rising time-average $\Dl$ and block-to-block price moves of about \eur{28} in both directions (Section~\ref{sec:occupation}). \citet{young1993} and \citet{kandori1993} study stochastic stability, the selection of states as the noise in a learning dynamic vanishes: the surviving states are those that need the most rare draws to leave and the fewest to reach, which in $2\times 2$ coordination games is the equilibrium with the larger basin of attraction under the noisy dynamic, not necessarily the one with the better static properties. % TODO-FACT: 06, 07 and FACTS §8 say "under the noisy dynamic" and this sentence matches them; the checker notes Kandori--Mailath--Rob select by basin size under the UNPERTURBED dynamic. One phrasing in all four places; FACTS to decide.
 This paper uses that idea to read the frozen outcome of each learner as the small-noise limit of its occupation measure, without claiming a formal result and, for the memory-1 learner, on the strength of a single low-noise point. \citet{fudenberg1986} state the folk theorem with discounting; its relevance here is negative, since a threat needs a strategy that conditions on history, and a memory-0 learner has none. The working definition of tacit collusion used here, that a firm holds back because it expects retaliation, follows \citet{tirole1988} and \citet{ivaldi2003}. \citet{watkins1992} prove convergence of Q-learning for a single agent in a stationary environment. \citet{busoniu2008} survey multi-agent reinforcement learning, where no general guarantee exists because each learner makes the others' environment non-stationary. A run that has stopped changing has therefore reached a rest point of the update on the actions it plays, not necessarily a fixed point of the Bellman operator on the full table. Section~\ref{sec:residual} measures that gap at one entry per temptation, the best unplayed deviation at each state of the converged cycle, with a one-step backup whose continuation value comes from the same table. \citet{oecd2017} and \citet{harrington2018} pose the policy question: what evidence would justify treating a price set by autonomous agents as collusion. This paper's contribution is narrower: in one simulated market the behavioural indicators do not separate a table that carries a continuation from one that carries a stale estimate, except where the punishment test records the state definition itself (a one-state table has nothing to react to), and a one-step test on the Q-table does, subject to the caveat that its continuation value may itself be stale. Because that test needs the table, it does not by itself answer the policy question for observed bid data.
